\section{Conclusion}
\label{sec:conclusion}

Deciding which evidence to request is a frequent, consequential and rule-governed decision that language-model agents are beginning to make. We asked whether such an agent should apply the rules itself or let code apply them to its verified judgments of the case. \casepath shows that execution can be made sound under incomplete facts and exact in what it defers, and our experiments show what that buys: request quality that depends far less on the model, no premature requests by construction, and errors that can be traced to a rule or a judgment. It does not buy accuracy for free. With a frontier model, letting the model apply the rules was as safe and, on imperfect rules, more accurate; execution is only as good as its rules and its policy for unstated facts. Agents that act in regulated processes should therefore execute verified rules over judgments they record, and evaluations of such agents should span models, runs and the facts a case leaves open.
